
\documentclass[9pt,a4paper]{extarticle}

\usepackage[a4paper,margin=0.45in]{geometry}
\usepackage{enumitem}
\usepackage{lmodern}
\usepackage[T1]{fontenc}

\pagestyle{empty}
\setlength{\parindent}{0pt}
\setlength{\parskip}{1pt}

\begin{document}

\begin{minipage}[t][0.47\textheight][t]{\textwidth}

\begin{center}
{\large\bfseries GRAPHIC ERA DEEMED TO BE UNIVERSITY}\\[-1mm]
{\small Department of Computer Science and Engineering}\\[-1mm]
{\bfseries MID-SEMESTER PRACTICAL EXAMINATION}\\[-1mm]
{\small Operating Systems Laboratory
\quad | \quad \textbf{SET--3 }}\\[-1mm]
{\small \textbf{Maximum Marks: 20}}
\end{center}

\vspace{-2mm}
\textbf{Name:} \rule{4.2cm}{0.4pt}
\hfill
\textbf{Roll No.:} \rule{3.5cm}{0.4pt}

\vspace{1mm}
\hrule
\vspace{1mm}

\textbf{Section A: Programming Problems (2 $\times$ 6 = 12 Marks)}

\begin{enumerate}[
label=\textbf{P\arabic*.},
leftmargin=6mm,
itemsep=2pt,
topsep=2pt
]

\item \textbf{Multi-Process Computation and Synchronization [6]}

Write a C program using \texttt{fork()}, \texttt{waitpid()} and
\texttt{exit()} to create exactly three children. Given an
integer array, Child 1 computes the sum of prime numbers,
Child 2 computes the sum of perfect squares, and Child 3
computes the sum of numbers divisible by both 3 and 5.
Each child displays its PID, PPID and result. The parent
must collect all three termination statuses using
\texttt{waitpid()} before displaying its completion message.

\item \textbf{Bidirectional IPC Using Two Pipes [6]}

Write a C program using two unnamed pipes. The parent sends
an integer array to the child. The child determines the
\textbf{second-largest distinct element} and the
\textbf{sum of even elements}, then returns both results
to the parent. Handle arrays with fewer than two distinct
elements and close all unused pipe descriptors.

\end{enumerate}

\vspace{1mm}
\textbf{Section B: Short Answer Questions (4 $\times$ 2 = 8 Marks)}

\begin{enumerate}[
label=\textbf{Q\arabic*.},
leftmargin=6mm,
itemsep=2pt,
topsep=2pt
]

\item For \texttt{fork(); fork(); fork();}, determine
the total number of processes and newly created children.

\item What does \texttt{waitpid(pid, \&status, WNOHANG)}
return if the specified child is still running?

\item Explain the behavior of \texttt{read()} on an empty
pipe when a writer is open versus when all writers are closed.

\item How can incorrect handling of pipe descriptors cause
indefinite blocking? Give one example.

\end{enumerate}

\vfill
\hrule
\begin{center}
{\scriptsize\textbf{End of SET--3}
\quad | \quad All programs must be implemented in C on Linux.}
\end{center}

\end{minipage}

\vspace{1cm}
\begin{minipage}[t][0.47\textheight][t]{\textwidth}

\begin{center}
{\large\bfseries GRAPHIC ERA DEEMED TO BE UNIVERSITY}\\[-1mm]
{\small Department of Computer Science and Engineering}\\[-1mm]
{\bfseries MID-SEMESTER PRACTICAL EXAMINATION}\\[-1mm]
{\small Operating Systems Laboratory
\quad | \quad \textbf{SET--3 }}\\[-1mm]
{\small \textbf{Maximum Marks: 20}}
\end{center}

\vspace{-2mm}
\textbf{Name:} \rule{4.2cm}{0.4pt}
\hfill
\textbf{Roll No.:} \rule{3.5cm}{0.4pt}

\vspace{1mm}
\hrule
\vspace{1mm}

\textbf{Section A: Programming Problems (2 $\times$ 6 = 12 Marks)}

\begin{enumerate}[
label=\textbf{P\arabic*.},
leftmargin=6mm,
itemsep=2pt,
topsep=2pt
]

\item \textbf{Multi-Process Computation and Synchronization [6]}

Write a C program using \texttt{fork()}, \texttt{waitpid()} and
\texttt{exit()} to create exactly three children. Given an
integer array, Child 1 computes the sum of prime numbers,
Child 2 computes the sum of perfect squares, and Child 3
computes the sum of numbers divisible by both 3 and 5.
Each child displays its PID, PPID and result. The parent
must collect all three termination statuses using
\texttt{waitpid()} before displaying its completion message.

\item \textbf{Bidirectional IPC Using Two Pipes [6]}

Write a C program using two unnamed pipes. The parent sends
an integer array to the child. The child determines the
\textbf{second-largest distinct element} and the
\textbf{sum of even elements}, then returns both results
to the parent. Handle arrays with fewer than two distinct
elements and close all unused pipe descriptors.

\end{enumerate}

\vspace{1mm}
\textbf{Section B: Short Answer Questions (4 $\times$ 2 = 8 Marks)}

\begin{enumerate}[
label=\textbf{Q\arabic*.},
leftmargin=6mm,
itemsep=2pt,
topsep=2pt
]

\item For \texttt{fork(); fork(); fork();}, determine
the total number of processes and newly created children.

\item What does \texttt{waitpid(pid, \&status, WNOHANG)}
return if the specified child is still running?

\item Explain the behavior of \texttt{read()} on an empty
pipe when a writer is open versus when all writers are closed.

\item How can incorrect handling of pipe descriptors cause
indefinite blocking? Give one example.

\end{enumerate}

\vfill
\hrule
\begin{center}
{\scriptsize\textbf{End of SET--3}
\quad | \quad All programs must be implemented in C on Linux.}
\end{center}

\end{minipage}
\end{document}
